\documentclass[11pt,spanish]{article}
\usepackage{graphicx}
\usepackage[spanish]{babel}
\usepackage[T1]{fontenc}
\usepackage{amsmath}
\usepackage{amsthm}
\usepackage{amsfonts}
\usepackage{amssymb}
\usepackage[dvips]{epsfig}
\usepackage{indentfirst}
\usepackage{fancyhdr}
\usepackage{color}
\usepackage{framed}
\usepackage{xcolor}
\usepackage{tikz}
\usetikzlibrary{babel,shapes.geometric}


\usepackage[top=2cm,left=2cm,right=2cm,bottom=2cm]{geometry}
\definecolor{rosa}{RGB}{236,64,122}
\definecolor{rosaoscuro}{RGB}{194,24,91}

\pagestyle{empty}

% \tarea{fila}{semana inicio}{semana fin}{texto}
\newcommand{\tarea}[4]{%
  \node[anchor=east, text width=7cm, align=left]
       at (-0.2,#1-0.5) {\textbf{#1}\quad #4};
  \filldraw[fill=rosa, draw=black]
       ({2*#2-2},{#1-0.8}) rectangle ({2*#3},{#1-0.2});
}
\begin{document}

\noindent
\begin{minipage}[c]{0.15\textwidth}
  \includegraphics[width=1.5cm]{uchile_full}
\end{minipage}%
\begin{minipage}[c]{0.65\textwidth}
  \centering
  {\Large\textbf{Computación y Programación}}\\[4mm]
  {\Large\textbf{Proyecto Final, año 2026}}
\end{minipage}%
\begin{minipage}[c]{0.15\textwidth}
  \raggedleft
  \includegraphics[width=3cm]{logo_departamento_fisica.pdf}
\end{minipage}

\vspace{10mm}
\section{Información general}

\begin{tabular}{lcl}
  T\'itulo del proyecto&:& \\
  Fecha de entrega del proyecto &:& viernes 6 de noviembre de 2026\\
  Área &:& Programación Científica \\
\end{tabular}


\vspace{30mm}

\section{Participante(s)}

Máximo de tres participantes.\\

\begin{tabular}{lcl}
  Nombre &:& Aranzazu Begoña Aretxabala Palacios \\
  Correo electrónico &:&aranzazu.aretxabala@ug.uchile.cl \\
  Carrera &:& Licenciatura en Física\\
\end{tabular}

\vspace{20mm}

\begin{tabular}{lcl}
  Nombre &:& \\
  Correo electrónico &:& \\
  Carrera &:& \\
\end{tabular}

\vspace{20mm}

\begin{tabular}{lcl}
  Nombre &:& \\
  Correo electrónico &:& \\
  Carrera &:& \\
\end{tabular}

\newpage
\section{Resumen}

\begin{framed}
  \noindent Debe ser claro e informativo. Describa los principales
  problemas que pretende atacar, incluyendo objetivos, metodología y
  resultados esperados. Un buen resumen facilita una descripción
  adecuada y entendimiento de lo que intenta lograr. El largo máximo
  de esta sección es de esta página en tamaño de letra 11~pt.
\end{framed}
%aqui se escribe 
\textbf{principales problemas a resolver:} 
\begin{itemize}
\item \textbf{comparar precios a mano es lento y poco eficaz:} Para encontrar la opción mas rentable hay que buscar el producto por varias tiendas e ir anotando el precios. El programa automatiza la búsqueda y reune todo en una tabla.
\item \textbf{Información dispersa y los precios no tienen un formato en común:}Cada proveedor coloca nombres, unidades, capacidades y precio de manera distinta. Este programa extrae esos datos y los normaliza para que se pueda comprar de manera justa. Además parte del trabajo es identificar cuando dos publicaciones corresponden al mismo material (con misma capacidad y tipo) para no comprar productos distintos.
\item \textbf{Precios cambian y la información se va desactualizando:} Una comparación hecha a mano solo sirve por ese día que se realiza la búsqueda. Con este proyecto se puede volver a ejecutar cuando se necesita y se puede obtener los precios vigentes e incluso se pude guardar un historial para ver como va variando el precio.
\item \textbf{Presupuesto de estudiantes y laboratorios universitarios son limitados:} Las pequeñas diferencias de precios cuando se compran muchas unidades. Asi que encontrar la opción mas económica por producto ayuda a gastar mejor.   
\end{itemize}
\textbf{ items faltantes:}
\begin{itemize}
\item \textbf{objetivos:} El objetivo principal de este proyecto es desarrollar una herramienta en python que extraiga y compare, de manera automática, mediante \textit{" web scraping"} los precios y materiales de laboratorio en distintos proveedores para identificar la opción más conveniente.
\item \textbf{Metodología:} Para este proyecto se usará \textit{Python} para extraer mediante \texttt{requests} y \texttt{BeautifulSoup} los precios de los materiales de laboratorio desde sitios proveedores chilenos. Los datos se comparan con \texttt{pandas}, se almacenan en CSV y se visualizan con \texttt{matplotlib} para visualizar la opción más conveniente por producto.
\item \textbf{resultados esperados:} El principal resultado que queremos lograr es crear un programa funcional en \textit{Python} que extraiga el nombre, precio y disponibilidad de materiales de laboratorio desde los sitios web de proveedores chilenos seleccionados para poder obtener el producto deseado a un precio accesible. Además se espera que los gráficos generados con \texttt{matplotlib} se muestren de manera clara las diferencias de precios. Tambien queremos lograr una herramienta reutilizable que se pueda ejecutar de nuevo para actualizar los precios o ampliar la lista de productos
\end{itemize}
\newpage

\section{Descripción del proyecto, marco teórico y revisión de la literatura}

\begin{framed}
  \noindent Esta sección debe incluir una presentación general del
  problema a desarrollar, con un breve marco teórico. Describa los
  aspectos novedosos que intentaron implementar y presente una
  descripción del estado del arte en el tópico del proyecto. El largo
  máximo de esta sección es de tres páginas en tamaño de letra 11~pt,
  incluidas referencias.
\end{framed}
% Aqui se escribe
Aqui se debe refenciar los papers usados ........
\newpage

\section{Objetivos}

\begin{framed}
  \noindent Especifique claramente sus objetivos principales y
  objetivos secundarios. El largo máximo de esta sección es de esta
  página, en tamaño de letra 11~pt.
\end{framed}
\subsection*{Objetivo Principal:} Es desarrollar una herramienta en \textit{Python} que extraiga y compare, mediante \textit{ web scraping}, los precios de materiales de laboratorio ofrecidos por diferentes proveedores chilenos, con el propósito de identificar la opción más económica para cada producto.
\subsection*{Objetivos secundarios:} 
\begin{itemize}
\item Seleccionar los proveedores chilenos y verificar que sus sitios web permitan el \textit{scraping} según su archivo \texttt{robots.txt} y sus términos de uso.
\item Programar la extracción automática del nombre, precio y la disponibilidad de cada producto, usando las bibliotecas de \texttt{requests} y \texttt{BeautifulSoup}.
\item Limpiar y normalizar los datos obtenidos con \texttt{pandas} y almacenarlos en un archivo CSV. 
\item Aplicar técnicas de \textit{scraping} responsable, espaciando las peticiones para no sobrecargar los servidores de los proveedores.
\item Visualizar los resultados de la búsqueda mediante tablas y gráficos elaborados con \texttt{matplotlib}
\item Documentar el código y los resultados, de modo que el programa se pueda reutilizar para actualizar los precios y ampliar la lista de los productos.
\end{itemize} 
\newpage

\section{Metodología}

\begin{framed}
  \noindent Describa los métodos que planea utilizar para lograr los
  objetivos planteados. Por ejemplo: técnicas numéricas, módulos, etc.
  Justifique sus elecciones. El largo máximo de esta sección es dos páginas, en
  tamaño de letra 11~pt.
\end{framed}
%aqui se escribe
\newpage

\section{Plan de trabajo}

\begin{framed}
  \noindent En la base de sus objetivos, indique las etapas y describa
  las actividades que realizaron. El largo máximo de esta sección es de
  esta página, en tamaño de letra 11~pt. Si es necesario, incorpore una
  carta Gantt.
\end{framed}
El proyecto se desarrollará en seis semanas (entre el 2 de noviembre de 2026 y 11 de diciembre de 2026), y se organiza en las siguientes etapas, asociados a los objetivos planteados con anterioridad:
\begin{itemize}
\item \textbf{Redacción y entrega de propuesta definitiva (semana 1):} Se redactará la propuesta definitiva y se entregará el 6 de noviembre de 2026.
\item \textbf{selección de proveedores y revisión de \texttt{robots.txt} (semanas 1--2):} Se eligirán de dos a cinco proveedores de materiales de laboratorios chilenos y se hará una lista acotada de productos a comprar. Además se revisará el archivo \texttt{robots.txt} y la estructura HTML de cada sitio para hacer viable el \textit{scraping}.
\item \textbf{Programación de los \textit{scrspers} (semanas 2--3):} Se escribirán en \textit{Python} con \texttt{requests} y \texttt{BeautifulSoup}, los programas que van a extraer el nombre, precio y disponibilidad de cada producto. Se presentará el avance del proyecto el 20 de noviembre de 2026.
\item \textbf{Limpieza y normalización de datos (semanas 3--4):} Con \texttt{pandas} se unificarán los precios, nombres y unidades. Además se guardarán en un archivo CSV.
\item \textbf{Comparación de precios y gráficos (semanas 4--5):} Se compararán los precios entre los proveedores seleccionados anteriormente y se selecionará la opción más económica.También se elaborarán tablas y gráficos usando \texttt{matplotlib}.
\item \textbf{Pruebas, corrección, documentación del código y presentación final (semanas 4--5):} Se probará el programa, se corregirán los errores y se documentará el código en un informe para finalmente presentarlo el 11 de diciembre de 2026.
\end{itemize}
\begin{figure}[h]
\centering
\resizebox{0.9\textwidth}{!}{%
\begin{tikzpicture}[x=1cm, y=-0.9cm, font=\small]

  % Cuadrícula: verticales punteadas, horizontales sólidas
  \foreach \x in {0,2,...,12} \draw[gray, dashed] (\x,0) -- (\x,9);
  \foreach \y in {0,...,9}    \draw (0,\y) -- (12,\y);

  % Encabezado de semanas
  \foreach \w/\f in {1/{2--6 nov}, 2/{9--13 nov}, 3/{16--20 nov},
                     4/{23--27 nov}, 5/{30 nov--4 dic}, 6/{7--11 dic}}
    \node[align=center] at ({2*\w-1},-0.7) {\textbf{\w}\\[-2pt]\scriptsize \f};

  \node[ font=\bfseries\large] at (-3.8,-0.7) {Descripción de las tareas};

  % Tareas
  \tarea{1}{1}{1}{Redacción y entrega de la propuesta definitiva}
  \tarea{2}{1}{2}{Selección de proveedores y revisión de \texttt{robots.txt}}
  \tarea{3}{2}{3}{Programación de los \textit{scrapers}}
  \tarea{4}{3}{4}{Limpieza y normalización de datos}
  \tarea{5}{4}{5}{Comparación de precios y gráficos}
  \tarea{6}{5}{5}{Pruebas, corrección y documentación del código}
  \tarea{7}{4}{5}{Redacción del informe final}
  \tarea{8}{5}{6}{Preparación de la presentación}

  % Hitos del curso
  \node[anchor=east, text width=7cm, align=left] at (-0.2,8.5) {\textbf{9}\quad Hitos del curso};
  \foreach \x/\t in {2/06/11, 6/20/11, 10/07/12, 11.5/10--11/12} {
    \node[diamond, draw=black, fill=rosaoscuro, inner sep=2.5pt] at (\x,8.5) {};
    \node[below, font=\scriptsize] at (\x,9) {\t};
  }

\end{tikzpicture}%
}
\caption{Carta Gantt del proyecto (tiempo en semanas).}
\end{figure}
\newpage

\section{Información adicional}

\begin{framed}
  \noindent Indique información extra que usted considere relevante,
  no incluida en las secciones anteriores, que puedan ayudar a la
  revisión de este proyecto. El largo máximo de esta sección es de
  esta página, en tamaño de letra 11~pt.
\end{framed}
% aqui se escribe
\end{document}
